\Needspace{9\baselineskip}
\section{Annulation de la fonction complétée et fermeture spectrale sur la droite critique}
La représentation conventionnelle des zéros de Riemann commence généralement trop tard. On présente la fonction zêta, on la prolonge analytiquement et l’on demande en quels points elle s’annule. Le lecteur voit une fonction complexe et une collection de points mystérieux répartis sur une droite.\par
La lecture HMT change l’ordre de la question. Elle ne commence pas par regarder où l’expression finale s’annule. Elle commence par reconstruire l’état que cette expression a contracté : les horloges des puissances premières, le canal gamma, le canal polaire, les deux feuillets, leur orientation, la retenue, la mémoire, la frontière, le terme terminal et la connexion nonadique qui les transporte conjointement.\par
Ce n’est qu’ensuite que la fonction complétée apparaît et que la question de ses zéros est posée.\par
La différence est radicale : un zéro cesse d’être une absence et devient un événement de fermeture.\par
\Needspace{7\baselineskip}
\subsection{Annulation scalaire et persistance de l’état enrichi}
Si\par
\[
\xi(\rho)=0,
\]
le zéro est véritable. Il n’est ni approximatif, ni fictif, ni décoratif. La fonction complétée s’annule exactement en \(\rho\).\par
Mais ce qui s’annule est une expression scalaire. Le point \(\rho\) ne disparaît pas, sa multiplicité ne disparaît pas, la structure première qui participe à la formule explicite ne disparaît pas, et l’état enrichi précédant l’expression ne disparaît pas.\par
C’est une distinction élémentaire et pourtant décisive :\par
\[
\text{valeur publiée égale à zéro}
\quad\not\Rightarrow\quad
\text{état générateur inexistant}.
\]
Une représentation utile est celle d’une frange sombre produite par interférence. L’obscurité ne signifie pas qu’aucune onde n’est arrivée. Elle signifie exactement le contraire : des contributions sont arrivées dont la relation de phase est si précise que leur somme visible s’annule.\par
Les zéros non triviaux peuvent être lus de cette manière : comme des nœuds d’interférence arithmétique globale.\par
Ce ne sont pas les nœuds d’une seule suite de nombres premiers, et ils n’appartiennent pas individuellement à un nombre premier. Dans l’expression complète interviennent simultanément :\par
\begin{itemize}
\item les horloges des puissances premières,
\end{itemize}
\[
\ell_{p,k}=k\log p;
\]
\begin{itemize}
\item le canal archimédien gamma ;
\item le canal polaire, qui conserve la normalisation et retire les pôles et les zéros triviaux ;
\item les deux feuillets spéculaires ;
\item la frontière que la projection visible ne montre pas.
\end{itemize}
Un zéro est une annulation exacte de cette expression conjointe. Son existence révèle une coordination, non un vide.\par
Les zéros constituent donc des \textbf{nœuds spectraux d’annulation arithmétique} : l’annulation de \(\xi\) exprime une coordination globale exacte entre les contributions de l’expression conjointe.\par
\Needspace{7\baselineskip}
\subsection{Fixation de la droite critique par l’involution fonctionnelle \(s\mapsto1-\bar s\)}
La fonction complétée conserve les symétries\par
\[
\rho\longmapsto1-\rho,
\qquad
\rho\longmapsto\overline\rho.
\]
Leur combinaison fait apparaître l’involution\par
\[
\rho\longmapsto\rho^\sharp
:=
1-\overline\rho.
\]
Si\par
\[
\rho=\beta+i\gamma,
\]
alors\par
\[
\rho^\sharp
=
1-\beta+i\gamma.
\]
La réflexion conserve la hauteur \(\gamma\), mais échange \(\beta\) et \(1-\beta\). Son ensemble de points fixes est :\par
\[
\rho=\rho^\sharp
\quad\Longleftrightarrow\quad
\beta=1-\beta
\quad\Longleftrightarrow\quad
\beta=\frac12.
\]
Por tanto,\par
\[
\operatorname{Fix}(\sharp)
=
\left\{
\frac12+i\gamma:\gamma\in\mathbb R
\right\}.
\]
La droite critique n’est pas simplement « la moitié » par préférence géométrique. Elle est la couture où les deux feuillets fonctionnels cessent d’apparaître comme des positions distinctes et deviennent une seule position autoduale.\par
Hors de la droite, un zéro est lié à une orbite pouvant compter jusqu’à quatre points :\par
\[
\rho,\qquad
1-\rho,\qquad
\overline\rho,\qquad
1-\overline\rho.
\]
Sur la droite critique, la réflexion fonctionnelle conjuguée n’envoie plus le zéro vers une autre position latérale :\par
\[
1-\overline\rho=\rho.
\]
Le zéro se referme sur lui-même.\par
C’est le premier sens précis du cercle de fermeture : il ne s’agit pas seulement de deux moitiés dont la somme vaut un, mais d’une orbite spéculaire qui rencontre un point fixe.\par
Cela ne démontre toutefois pas encore l’hypothèse de Riemann. L’égalité\par
\[
\frac12=1-\frac12
\]
identifie l’axe ; elle n’impose pas à elle seule que tous les zéros s’y trouvent.\par
La contrainte provient de la positivité.\par
\Needspace{7\baselineskip}
\subsection{Quadraticité du produit croisé sous fermeture de l’involution fonctionnelle}
La forme spectrale complète peut s’écrire\par
\[
\mathscr I_{\mathrm{GW}}(f,g)
=
\sum_{\rho\in\mathscr Z}
\Phi_f(\rho)\,
\overline{\Phi_g(1-\overline\rho)},
\]
où \(\mathscr Z\) est le multiensemble des zéros non triviaux, avec conservation de leurs multiplicités.\par
Cette formule détermine comment les zéros sont appariés : l’évaluation située en \(\rho\) se combine avec celle située en son reflet \(\rho^\sharp=1-\overline\rho\).\par
Si le zéro est hors de la droite, l’expression apparie deux points différents :\par
\[
\Phi_f(\rho)\,
\overline{\Phi_f(\rho^\sharp)}.
\]
C’est un produit croisé. Ce n’est pas nécessairement un module au carré et, sur une classe de fonctions tests suffisamment riche, il peut produire des directions de signe négatif.\par
Si le zéro est sur la droite critique, alors\par
\[
\rho^\sharp=\rho
\]
et le même terme devient\par
\[
|\Phi_f(\rho)|^2.
\]
La forme complète devient :\par
\[
\mathscr I_{\mathrm{GW}}(f,f)
=
\sum_{\rho\in\mathscr Z}
|\Phi_f(\rho)|^2.
\]
Le miroir n’apparie plus deux positions décalées. Chaque point est apparié avec lui-même et l’expression devient une somme de carrés.\par
Telle est la lecture conceptuelle de la démonstration : HMT construit en amont la forme de Weil comme une énergie de frontière,\par
\[
\mathscr I_{\mathrm{GW}}
=
E_R^*E_R
=
\mathscr L_R^*\mathscr L_R
\succeq0.
\]
Elle n’est pas déclarée positive parce que les zéros seraient sur la droite. Elle est construite comme carré avant toute utilisation des zéros.\par
Ensuite, le critère réversible de Weil affirme que cette positivité sur toute la classe séparante de fonctions tests équivaut au fait que tous les zéros non triviaux satisfont\par
\[
\operatorname{Re}\rho=\frac12.
\]
La positivité n’« attire » donc pas dynamiquement les zéros vers le centre. Les zéros ne se déplacent pas depuis une position initiale. Ce qui se produit est plus fort : toute configuration hors de l’axe fixe produirait une direction négative détectable, tandis que le complément HMT est matériellement une norme au carré et ne peut en acquérir une.\par
La droite critique est l’unique lieu compatible avec les deux expressions.\par
\Needspace{7\baselineskip}
\subsection{Caractérisation spectrale des zéros}
Dans cette lecture, un zéro non trivial peut être décrit simultanément de quatre manières compatibles.\par
C’est un \textbf{véritable zéro analytique}, puisque\par
\[
\xi(\rho)=0.
\]
C’est un \textbf{point fixe du miroir}, puisque\par
\[
\rho=1-\overline\rho.
\]
C’est un \textbf{atome spectral positif}, puisqu’il apparaît dans la forme de Weil comme\par
\[
|\Phi_f(\rho)|^2.
\]
Et c’est un \textbf{nœud global des horloges premières et du bord archimédien}, car il n’appartient pas à un canal isolé : il émerge de l’expression coordonnée des termes premier, gamma et polaire.\par
Le zéro est nul comme amplitude et positif comme localisation spectrale. Il n’y a aucune contradiction. Ce sont deux lecteurs différents.\par
La fonction vaut zéro en \(\rho\), mais \(\rho\) continue d’apparaître, avec toute sa multiplicité, dans la mesure spectrale positive. Le zéro n’a pas été effacé parce qu’il est nul. Au contraire : il devient l’un des lieux qui soutiennent la somme positive.\par
C’est probablement le changement interprétatif le plus important :\par
\begin{quote}
Le zéro n’est pas le lieu où la structure disparaît. C’est le lieu où la structure complète atteint une annulation visible et laisse une marque spectrale.\par
\end{quote}
\Needspace{7\baselineskip}
\subsection{Identification de la hauteur d’un \(0\) non trivial à une fréquence spectrale réelle}
La normalisation utilisée dans le critère de Weil est\par
\[
\Phi_f(s)
=
\widehat f\!\left(
i\left(s-\frac12\right)
\right).
\]
Si\par
\[
\rho=\frac12+i\gamma,
\]
alors\par
\[
i\left(\rho-\frac12\right)
=
-\gamma
\]
et, par conséquent,\par
\[
\Phi_f(\rho)
=
\widehat f(-\gamma).
\]
La hauteur \(\gamma\) devient littéralement une fréquence réelle de Fourier.\par
La droite critique peut donc être lue comme une droite de fréquences : la coordonnée réelle \(\tfrac12\) fixe l’équilibre spéculaire et la coordonnée \(\gamma\) enregistre la fréquence du nœud.\par
Cela explique pourquoi le langage spectral est approprié. Les zéros ne sont pas seulement des coordonnées complexes ; ce sont des fréquences auxquelles l’expression arithmétique complète manifeste sa condition nodale.\par
Ces hauteurs admettent l’interprétation de fréquences nodales d’un système global : des points où une amplitude observable s’annule tandis que la structure produisant cette annulation demeure active.\par
Mais il faut protéger cette phrase contre une exagération. Il n’en découle pas automatiquement que chaque \(\gamma\) soit une valeur propre énergétique d’un Hamiltonien physique concret. Il faudrait pour cela construire ce Hamiltonien, son domaine et la correspondance spectrale. Ce qui est démontré est la lecture de Fourier–Mellin, la forme positive complète et la localisation des zéros.\par
Il s’agit rigoureusement de fréquences spectrales. Elles ne deviennent pas, sans un autre entrelaceur, des niveaux physiques d’énergie.\par
\Needspace{7\baselineskip}
\subsection{Paramétrisation circulaire conforme de la droite critique}
Votre intuition du cercle n’est pas seulement métaphorique. Il existe une transformation exacte de Möbius :\par
\[
\tau(s)
=
\frac{s-1}{s}.
\]
Se cumple:\par
\[
\operatorname{Re}s=\frac12
\quad\Longleftrightarrow\quad
|\tau(s)|=1.
\]
La droite critique compactifiée par son point à l’infini se transforme donc en cercle unité.\par
Si\par
\[
s=\frac12+it,
\]
alors\par
\[
\tau(s)\in S^1.
\]
Al escribir\par
\[
\tau=e^{i\theta},
\]
l’inverse est\par
\[
s=\frac1{1-\tau}
=
\frac12+\frac{i}{2}\cot\frac{\theta}{2}.
\]
La hauteur \(t\) de la droite devient une position angulaire sur le cercle. Lorsque \(t\) parcourt toute la droite, la carte parcourt le cercle compactifié. Les deux extrémités \(t\to-\infty\) et \(t\to+\infty\) se rejoignent au point \(\tau=1\), bien qu’elles le fassent selon des orientations opposées.\par
Le miroir prend lui aussi une forme circulaire :\par
\[
\tau(1-\overline s)
=
\frac1{\overline{\tau(s)}}.
\]
L’involution spectrale devient une réflexion par rapport au cercle unité, dont l’ensemble fixe est précisément le cercle.\par
L’expression « cercle de fermeture » possède donc un contenu exact :\par
\begin{itemize}
\item la droite critique compactifiée est un cercle ;
\item le miroir fonctionnel devient une réflexion par rapport à ce cercle ;
\item les zéros non triviaux sont des points marqués sur cette frontière ;
\item chaque point est fixé par la réflexion qui échangeait auparavant les deux feuillets.
\end{itemize}
Le cercle n’engendre pas les zéros et ne détermine pas leurs hauteurs. Toute la droite est transformée en cercle, tandis que l’équation\par
\[
\xi(\rho)=0
\]
sélectionne ses points spectraux.\par
Cela ne signifie pas non plus que les zéros soient équidistants ou qu’il y en ait neuf par tour. Le cercle est la géométrie compactifiée du lieu critique ; la distribution des zéros sur celui-ci conserve son arithmétique propre.\par
\Needspace{7\baselineskip}
\subsection{Cercle nonadique et holonomie de phase}
Il existe un second cercle, antérieur au cercle analytique : la base cyclique des neuf phases.\par
La connexion nonadique ne se compose ni de neuf filtres ni de neuf canaux indépendants. C’est une unique holonomie :\par
\[
\Gamma_9
=
\operatorname{Hol}_{C_9}
(\nabla_{\mathrm{TPK}}).
\]
Sa loi de retour est\par
\[
g(K+9)=g(K),
\]
mais simultanément\par
\[
q_{\mathrm{mem}}
(\Gamma_9\widehat x)
=
q_{\mathrm{mem}}(\widehat x)+1
\]
et, en général,\par
\[
\Gamma_9\widehat x
\neq
\widehat x.
\]
La phase revient ; l’état n’est pas réinitialisé.\par
Vu depuis la base, un cercle a été parcouru et l’on est revenu à la même phase. Vu depuis l’état enrichi, un tour de mémoire s’est accumulé. L’image la plus fidèle pour le lecteur est :\par
\begin{itemize}
\item un cercle dans la projection ;
\item une hélice dans l’état complet.
\end{itemize}
Vue d’en haut, l’hélice paraît se refermer. Vue de l’intérieur, chaque tour occupe une hauteur différente.\par
Ce qui est littéral n’est pas l’image de l’hélice, mais les identités mathématiques : retour de phase, incrément de mémoire et changement général de l’état.\par
Cette structure a été vérifiée sur 396 niveaux de raffinement par canal. Le certificat enregistre 43 tours nonadiques complets et 387 comparaisons \(K,K+9\) par canal : la phase coïncide et l’état change dans chacune d’elles. Ce n’est pas une manière poétique de parler de périodicité ; c’est une séparation effectivement réalisée entre retour visible et continuité généalogique.\par
\Needspace{7\baselineskip}
\subsection{Le miroir nonadique et le miroir de Riemann}
Votre intuition est correcte : dans les deux cas apparaît un miroir. Il convient toutefois de préciser exactement leur relation.\par
Le miroir interne nonadique échange les deux composantes latérales associées aux lectures \(\pi/e\), tout en conservant l’axe \(\varphi\) et l’agrégat correspondant. L’involution de Riemann est\par
\[
s\longmapsto1-\overline s.
\]
Ce ne sont pas littéralement la même application, car leurs domaines et codomaines sont distincts. Les égaler sans entrelaceur reviendrait à confondre l’architecture interne avec son expression analytique.\par
L’assertion correcte est plus intéressante : l’entrelaceur spectral–polaire exprime l’architecture bidirectionnelle du TPK comme l’involution fonctionnelle de la fonction complétée. Le miroir nonadique est antérieur ; le miroir de Riemann en est la réalisation dans le codomaine analytique.\par
Dans les deux cas existent :\par
\begin{itemize}
\item un couple de feuillets ;
\item un axe qui demeure fixe ;
\item une transformation qui échange les positions latérales ;
\item une couture où les deux expressions coïncident ;
\item une mémoire qui ne disparaît pas lors de l’identification.
\end{itemize}
La droite \(\Re s=\tfrac12\) est la couture visible de ce miroir dans le plan complexe.\par
\Needspace{7\baselineskip}
\subsection{Fermeture résiduelle et phase nonadique}
L’expression « les neuf sont de faux zéros » contient une intuition très précise si l’on distingue l’entier, la classe résiduelle et l’état complet.\par
Pour \(n>0\), la carte positive utilise\par
\[
\rho_9(n)
=
1+((n-1)\bmod9).
\]
De manière équivalente,\par
\[
\rho_9(n)
=
\begin{cases}
n\bmod9,
&
n\not\equiv0\pmod9,\\[1mm]
9,
&
n\equiv0\pmod9.
\end{cases}
\]
Par conséquent, lorsqu’un nombre appartient à la classe nulle modulo neuf, la carte positive n’exprime pas \(0\) : elle exprime \(9\).\par
Simultanément sont conservés\par
\[
\kappa_9(n)
=
\frac{n-\rho_9(n)}9
\]
et la reconstruction\par
\[
n
=
9\kappa_9(n)
+
\rho_9(n).
\]
Le neuf est ainsi le représentant positif de la classe résiduelle zéro. Il est nul dans le quotient, mais il n’est ni l’entier zéro ni un état vide.\par
Le neuf signifie : « le tour est achevé ». Le quotient indique combien de couronnes ou de tours se sont accumulés. Si tout est remplacé par le résidu \(0\), la congruence est conservée, mais la différence entre absence, fermeture et mémoire est perdue.\par
C’est le sens rigoureux de « faux zéro » :\par
\begin{quote}
Le neuf est un zéro pour l’image résiduelle, mais il serait faux de l’interpréter comme une annihilation de l’état. C’est un zéro exprimé avec mémoire de fermeture.\par
\end{quote}
Par exemple, \(27\), \(72\) et \(729\) possèdent une projection résiduelle nulle et une racine numérique positive égale à neuf. Cependant, ils ne constituent ni le même nombre ni la même généalogie. L'un peut provenir d'une réversion, un autre d'une évaluation différente et un autre d'une puissance. Le résidu ne reconstitue pas le parcours.\par
Cela établit une connexion profonde avec les zéros de Riemann, sans toutefois les identifier.\par
Un zéro de \(\xi\) est un véritable zéro analytique. Ce n'est pas le nombre neuf. Rien n'affirme ici que chaque zéro non trivial corresponde individuellement à une phase neuf particulière.\par
La relation structurelle est la suivante :\par
\[
\text{une publication peut s'annuler}
\quad\text{sans que sa généalogie s'annule}.
\]
Le neuf est une fausse absence arithmétique. Le zéro de zêta est une fausse absence ontologique. Tous deux sont de véritables zéros dans leur lecteur, mais aucun ne signifie qu'aucune structure ne subsiste en arrière-plan.\par
\Needspace{7\baselineskip}
\subsection{Distinction typée entre les zéros de \(\zeta\), de \(\xi\) et du déterminant spectral}
Pour que la force de la narration n'efface pas les types mathématiques, il convient de distinguer trois zéros.\par
\Needspace{5\baselineskip}
\subsubsection{\(0\) analytique comme annulation de la fonction dans le domaine spectral}
\[
\xi(\rho)=0.
\]
C'est un point spectral de la fonction complétée. Il conserve position, hauteur, multiplicité et symétries.\par
\Needspace{5\baselineskip}
\subsubsection{\(0\) terminal comme objet universel du morphisme de clôture}
Dans la télescopie HMT apparaît\par
\[
(A^NE_R)^*(A^NE_R)
=
q^{2N}
\mathcal B_{\mathrm{sp},R}^{\mathrm{str}}
\longrightarrow0,
\qquad
q=\frac59.
\]
C'est l'extinction d'un reste après le transfert de toute sa contribution vers une somme de carrés positifs.\par
\Needspace{5\baselineskip}
\subsubsection{Représentant positif de la classe résiduelle nulle : \(n\equiv0\pmod 9\) et \(\rho_9(n)=9\)}
\[
n\equiv0\pmod9,
\qquad
\rho_9(n)=9.
\]
C'est la classe modulaire nulle exprimée par un représentant positif, accompagnée de son quotient et de sa mémoire.\par
Il ne s'agit pas du même objet. Ils partagent cependant une grammaire :\par
\[
\text{annulation visible}
+
\text{conservation structurelle}.
\]
\Needspace{7\baselineskip}
\subsection{Épuisement de l'objet terminal sous la filtration cofinale}
La démonstration conserve à chaque profondeur une identité finie :\par
\[
\mathcal B_{\mathrm{sp},R}^{\mathrm{str}}
=
\sum_{n=0}^{N-1}
(DA^nE_R)^*(DA^nE_R)
+
(A^NE_R)^*(A^NE_R).
\]
Le dernier terme est le terminal. Il n'est pas supprimé pour obtenir une somme positive plus élégante. Il reste présent à chaque étape finie.\par
Ce n'est qu'ensuite que l'on démontre :\par
\[
(A^NE_R)^*(A^NE_R)
=
q^{2N}
\mathcal B_{\mathrm{sp},R}^{\mathrm{str}}
\longrightarrow0.
\]
Cela permet une autre lecture du zéro : le terminal tend vers zéro parce que son contenu a été exprimé et comptabilisé niveau après niveau. Chaque partie perdue par la contraction visible a été enregistrée dans une coordonnée de mémoire positive.\par
Le reste s'éteint, mais il n'a pas été effacé.\par
C'est comme vider un réservoir en transférant exactement son contenu vers une suite de récipients numérotés. À la fin, le réservoir initial est vide, mais rien n'a disparu : il existe un registre complet du transfert.\par
Ici, le zéro signifie \textbf{transfert achevé}.\par
\Needspace{7\baselineskip}
\subsection{Information généalogique du continu codée par les zéros critiques}
La droite critique\par
\[
\frac12+i\mathbb R
\]
est une copie translatée de la droite réelle. Dans HMT, cette droite n'est pas l'objet initial : elle est une expression archimédienne du continu généalogique.\par
Les zéros forment un ensemble discret au sein de cette expression continue. Ils ne perforent pas la droite et n'interrompent pas le continu. Ils constituent un diviseur spectral : des points isolés, chacun muni de sa multiplicité, sur une projection continue.\par
Cette coexistence du continu et du discret est essentielle. La droite fournit le codomaine continu ; la généalogie détermine les marques discrètes qui y apparaissent.\par
L'image que je proposerais au lecteur est celle d'une membrane continue traversée par des nœuds spectraux. Les nœuds ne rompent pas la membrane. Ils rendent visible sa structure interne.\par
En termes holographiques, le zéro est une marque de frontière produite par une organisation globale de l'intérieur. Il n'appartient pas à un nombre premier unique, de même qu'un pixel d'une image holographique n'appartient pas à un seul point de l'objet représenté. Chaque zéro répond à l'ensemble complet des relations premier–gamma–polaire.\par
Cela ne signifie pas que chaque zéro contienne à lui seul toute l'arithmétique sous une forme récupérable. Cela signifie que sa position ne peut être attribuée à une entrée locale isolée : c'est une réponse globale du système.\par
\Needspace{7\baselineskip}
\subsection{Correspondance entre quantification spectrale, fréquences réelles et réalisation physique de l'opérateur}
Les zéros s'interprètent comme des nœuds d'un transfert arithmétique global :\par
\begin{itemize}
\item la coordonnée \(\gamma\) agit comme fréquence ;
\item la droite critique est la surface d'équilibre spéculaire ;
\item la forme de Weil est l'énergie positive du complément ;
\item la connexion nonadique conserve la mémoire de chaque tour ;
\item le terminal mesure ce qui n'a pas encore été exprimé ;
\item le zéro visible signale une annulation parfaite de l'amplitude exprimée.
\end{itemize}
Ils ne s'identifient ni à des particules, ni à des trous, ni à des états du vide, ni à des masses, ni à des niveaux énergétiques matériels sans construire auparavant un opérateur physique supplémentaire.\par
L'interprétation physique valide est structurelle : nœud, fréquence, miroir, transfert, mémoire et clôture. L'ontologie physique concrète exigerait son propre transducteur.\par
\Needspace{7\baselineskip}
\subsection{Corollaires de localisation, de multiplicité et de symétrie de la caractérisation spectrale}
Avant la résolution, la question semblait être :\par
\begin{quote}
Pourquoi des points définis par une fonction complexe semblent-ils s'aligner mystérieusement sur une droite ?\par
\end{quote}
Après la résolution, la question peut se formuler autrement :\par
\begin{quote}
Comment la structure spectrale–polaire exprime-t-elle une énergie de frontière qui ne peut rester positive que lorsque chaque zéro s'identifie à son reflet ?\par
\end{quote}
La réponse est la suivante :\par
\begin{itemize}
\item les colonnes positive et négative proviennent du même état ;
\item l'état de frontière \(E_R\) conserve ce que la projection ne montre pas ;
\item la différence des deux matrices de Gram est
\end{itemize}
\[
E_R^*E_R;
\]
\begin{itemize}
\item la télescopie conserve le terminal jusqu'à démontrer son extinction ;
\item la forme de Weil complète coïncide avec ce carré ;
\item le critère réversible impose que chaque zéro soit un point fixe du miroir.
\end{itemize}
La droite n'est plus une coïncidence visuelle. Elle est l'ensemble fixe de l'involution compatible avec la positivité construite.\par
\Needspace{7\baselineskip}
\subsection{Hypothèses et codomaine exacts du théorème de localisation spectrale}
Cette lecture permet d'affirmer :\par
\[
\{\text{zéros non triviaux}\}
\subset
\left\{
\frac12+i\gamma:\gamma\in\mathbb R
\right\}.
\]
Et, au moyen de la carte de Möbius,\par
\[
\left\{
\frac12+i\gamma
\right\}
\longrightarrow
S^1.
\]
Par conséquent, tous les zéros non triviaux sont des points marqués du cercle critique.\par
On n'a pas construit ici de bijection individuelle :\par
\[
\{
\text{zéros non triviaux}
\}
\longleftrightarrow
\{
\text{tours nonadiques concrets}
\}.
\]
On n'affirme ni que le premier zéro soit le premier tour, ni qu'il existe neuf zéros par cycle, ni que leurs hauteurs soient engendrées par répétition d'une phase. La résolution démontre la localisation, la positivité, la clôture spéculaire et l'appartenance au cercle. Une quantification individuelle des hauteurs constituerait un développement supplémentaire et nécessiterait sa propre application.\par
Il ne faut pas non plus confondre le miroir interne \(\pi/e\), l'involution fonctionnelle de Riemann et une éventuelle symétrie physique. Leur relation relève d'une expression typée, non d'une identité nominale.\par
\Needspace{7\baselineskip}
\subsection{Théorème de clôture spectrale par factorisation positive de Guinand–Weil}
L'interprétation opératorielle se résume par les propriétés suivantes :\par
\begin{quote}
Les zéros non triviaux de Riemann sont des annulations exactes de l'amplitude analytique qui conservent la généalogie de leur construction. Chaque zéro détermine une fréquence globale, un marqueur spectral et un point autodual de l'involution fonctionnelle. La compactification de la droite critique produit le cercle fixe commun aux deux feuillets. La connexion nonadique distingue le retour de phase de la réinitialisation de l'état : la clôture circulaire conserve l'avancement de la mémoire. La classe nonadique \(9\equiv0\pmod 9\) représente la clôture visible accompagnée de l'information interne de l'état enrichi.\par
\end{quote}
L'annulation spectrale constitue ainsi une forme concentrée de clôture opératorielle et conserve l'objet généalogique sous-jacent.\par
Cette interprétation physico-mathématique accompagne les résultats démontrés sans s'y substituer et n'affirme aucune correspondance individuelle entre zéros et tours nonadiques tant que l'entrelaceur correspondant n'est pas construit.\par

\section{Forme de Weil, translation et quantification des hauteurs spectrales}
L'objet correct n'est pas un tour nu, mais une orbite hélicoïdale enrichie :\par
\[
\boxed{
\text{zéro non trivial}
\longleftrightarrow
\text{atome spectral pondéré}
\longleftrightarrow
\text{caractère hélicoïdal nonadique}.
}
\]
Cette formalisation intervient après le continu complet et ne redéfinit ni ne sépare ses cinq caractéristiques.\par
\Needspace{7\baselineskip}
\subsection{Construction de l'espace de Hilbert à partir de la forme de Weil}
Se toma\par
\[
\mathcal D=C_c^\infty(\mathbb R),
\qquad
\langle f,g\rangle_W
:=
\mathscr I_{\mathrm{GW}}(f,g).
\]
Dans la résolution HMT, cette forme ne naît pas de la liste des zéros. Elle se construit d'abord à partir des canaux premier–puissance, gamma et polaire, puis se factorise par la perte nonadique :\par
\[
\mathscr I_{\mathrm{GW}}
=
E^*E
=
\mathscr L^*\mathscr L
\succeq0.
\]
Se define\par
\[
\mathcal N_W
=
\{f\in\mathcal D:
\mathscr I_{\mathrm{GW}}(f,f)=0\},
\]
y completamos:\par
\[
\mathcal H_W
=
\overline{\mathcal D/\mathcal N_W}.
\]
Tel est l'espace positif de Weil construit du côté arithmétique. Les zéros n'ont pas encore été utilisés pour le définir.\par
Parallèlement, la télescopie fournit l'espace nonadique de perte\par
\[
\mathcal K_{\mathrm{loss}}
=
\overline{\operatorname{Ran}\mathscr L},
\]
et l'égalité des normes définit une unité canonique\par
\[
W_{\mathrm{loss}}:
\mathcal H_W
\longrightarrow
\mathcal K_{\mathrm{loss}},
\qquad
W_{\mathrm{loss}}[f]=\mathscr Lf.
\]
Nous ne plaçons pas après coup une copie artificielle de la forme de Weil dans le TPK : l'espace de perte de la preuve lui-même est une réalisation de son espace de Hilbert positif.\par
\Needspace{7\baselineskip}
\subsection{Opérateur de translation indépendant de la localisation préalable des zéros}
Sur le domaine global \(\mathcal D\), et non dans une fenêtre fixe, on considère\par
\[
(T_tf)(x)=f(x-t).
\]
L'invariance peut se vérifier avant de diagonaliser la formule sur les zéros. Avec la convention\par
\[
\widehat{T_tf}(z)
=
e^{-izt}\widehat f(z),
\]
se tiene\par
\[
h_{T_tf,T_tg}(z)=h_{f,g}(z),
\qquad
\check h_{T_tf,T_tg}(a)=\check h_{f,g}(a).
\]
Par conséquent, tous les termes premier, gamma et polaire restent invariants :\par
\[
\mathscr I_{\mathrm{GW}}(T_tf,T_tg)
=
\mathscr I_{\mathrm{GW}}(f,g).
\]
La translation conserve \(\mathcal N_W\) et induit un groupe unitaire :\par
\[
U_t[f]=[T_tf].
\]
Après la localisation critique déjà démontrée, sa continuité forte découle de\par
\[
\|U_t[f]-[f]\|_W^2
=
\sum_\gamma
m_\gamma
\left|e^{it\gamma}-1\right|^2
\left|\widehat f(-\gamma)\right|^2
\longrightarrow0.
\]
Le théorème de Stone produit alors un opérateur autoadjoint \(H_W\) :\par
\[
U_t=e^{itH_W}.
\]
Cet opérateur a été défini au moyen des translations et de la forme premier–gamma–polaire. Il n'a pas été défini en écrivant préalablement une matrice diagonale contenant les zéros.\par
\Needspace{7\baselineskip}
\subsection{Identification spectrale des hauteurs des zéros non triviaux}
Puisque l'hypothèse de Riemann est déjà démontrée, chaque zéro non trivial est de la forme\par
\[
\rho_\gamma=\frac12+i\gamma.
\]
Si \(\Gamma_\xi\) est l'ensemble des hauteurs distinctes et \(m_\gamma\) leur multiplicité, on introduit la mesure de comptage pondérée\par
\[
\mu_\xi
=
\sum_{\gamma\in\Gamma_\xi}
m_\gamma\,\delta_\gamma.
\]
La formule spectrale de Weil devient\par
\[
\mathscr I_{\mathrm{GW}}(f,g)
=
\sum_{\gamma\in\Gamma_\xi}
m_\gamma\,
\widehat f(-\gamma)
\overline{\widehat g(-\gamma)}.
\]
Por tanto,\par
\[
\mathcal A[f](\gamma)
=
\widehat f(-\gamma)
\]
définit une isométrie\par
\[
\mathcal A:
\mathcal H_W
\longrightarrow
L^2(\Gamma_\xi,\mu_\xi).
\]
En outre, aucune partie spectrale ne reste cachée hors de son image. L'image de \(\mathcal A\) est fermée et préservée par tous les multiplicateurs\par
\[
M_t a(\gamma)=e^{it\gamma}a(\gamma).
\]
La projection sur un atome individuel s'obtient par la moyenne spectrale\par
\[
P_\gamma
=
\operatorname*{s-lim}_{T\to\infty}
\frac1{2T}
\int_{-T}^{T}
e^{-it\gamma}M_t\,dt.
\]
La séparation des évaluations des transformées de fonctions tests garantit que, pour chaque \(\gamma\), il existe un \(f\) tel que\par
\[
\widehat f(-\gamma)\neq0.
\]
Par conséquent, \(P_\gamma\operatorname{Ran}\mathcal A\neq0\). Puisque l'image est réductrice, elle contient l'axe atomique correspondant. Tous les axes atomiques appartiennent à l'image et engendrent densément \(L^2(\mu_\xi)\). En conséquence,\par
\[
\boxed{
\mathcal H_W
\simeq
L^2(\Gamma_\xi,\mu_\xi)
}
\]
y\par
\[
\boxed{
\mathcal A H_W\mathcal A^{-1}
=
M_\gamma.
}
\]
Ainsi,\par
\[
\boxed{
\sigma(H_W)=\Gamma_\xi,
}
\]
sans hauteur supplémentaire, sans hauteur omise et sans composante continue.\par
C'est déjà une quantification spectrale individuelle : les hauteurs des zéros sont exactement le spectre ponctuel d'un opérateur défini à partir de la forme arithmétique positive.\par
La multiplicité exige une précision. La forme scalaire enregistre\par
\[
\mu_\xi(\{\gamma\})=m_\gamma.
\]
Autrement dit, elle conserve exactement le poids algébrique du zéro. Elle ne produit pas automatiquement \(m_\gamma\) vecteurs distinguables : si un zéro était multiple, l'évaluation au même point se répéterait avec le poids \(m_\gamma\). La correspondance canonique relie le diviseur des zéros aux atomes spectraux pondérés. Séparer intérieurement les copies exigerait une forme positive de jets ou une démonstration de simplicité.\par
\Needspace{7\baselineskip}
\subsection{Transformation conforme de la droite critique en cercle spectral avec conservation des hauteurs}
On utilise la normalisation\par
\[
\tau(s)=\frac{s-1}{s}.
\]
Alors\par
\[
\Re s=\frac12
\quad\Longleftrightarrow\quad
|\tau(s)|=1.
\]
Pour\par
\[
\rho_\gamma=\frac12+i\gamma
\]
se obtiene\par
\[
\tau_\gamma
=
\tau(\rho_\gamma)
=
\frac{-\frac12+i\gamma}{\frac12+i\gamma}
=
\frac{\gamma+\frac i2}{\gamma-\frac i2}
\in S^1.
\]
C'est exactement la transformée de Cayley de l'opérateur \(H_W\) :\par
\[
\mathcal C_W
=
\left(H_W+\frac i2I\right)
\left(H_W-\frac i2I\right)^{-1}.
\]
Sur l'atome \(\gamma\),\par
\[
\mathcal C_WP_\gamma
=
\tau_\gamma P_\gamma.
\]
La correspondance ne perd aucune information. Son inverse est\par
\[
\rho_\gamma=\frac1{1-\tau_\gamma}.
\]
Si\par
\[
\tau_\gamma=e^{i\theta_\gamma},
\]
alors\par
\[
\boxed{
\gamma
=
\frac12\cot\frac{\theta_\gamma}{2}.
}
\]
Par conséquent, la hauteur complète est contenue dans l'angle. La compactification de la droite conserve la localisation de chaque zéro.\par
De plus,\par
\[
\tau_{-\gamma}
=
\overline{\tau_\gamma}
=
\tau_\gamma^{-1}.
\]
Les zéros conjugués correspondent à des orientations inverses du même cercle. C'est le miroir converti en une identité angulaire exacte.\par
La correspondance démontrée est\par
\[
\boxed{
\sum_\rho m_\rho[\rho]
\longleftrightarrow
\sum_\gamma m_\gamma
[\tau_\gamma,P_\gamma].
}
\]
Elle ne dépend pas d'un classement préalable des zéros par hauteur.\par
\Needspace{7\baselineskip}
\subsection{Monodromie de chaque \(0\) sur le cercle spectral}
Transportons maintenant l'opérateur circulaire sur l'espace nonadique :\par
\[
\mathcal C_{\mathrm{loss}}
=
W_{\mathrm{loss}}\,
\mathcal C_W\,
W_{\mathrm{loss}}^{-1}.
\]
Dans la télescopie HMT,\par
\[
q=\frac59
\]
est la contraction correspondant à un tour complet. On définit la réalisation spectrale de ce tour :\par
\[
\mathcal M_{\mathrm{spec}}
=
q\,\mathcal C_{\mathrm{loss}}.
\]
Si \(v_\gamma\) appartient à l'atome spectral \(P_\gamma\), alors\par
\[
\boxed{
\mathcal M_{\mathrm{spec}}^n v_\gamma
=
q^n\tau_\gamma^n v_\gamma.
}
\]
C'est l'orbite hélicoïdale du zéro.\par
Son rayon est\par
\[
\|\mathcal M_{\mathrm{spec}}^n v_\gamma\|
=
q^n\|v_\gamma\|,
\]
tandis que son angle accumulé est\par
\[
n\theta_\gamma.
\]
La partie radiale conserve l'extinction terminale ; la partie angulaire conserve l'identité individuelle du zéro. Et toutes deux satisfont la télescopie :\par
\[
\|v_\gamma\|^2
=
(1-q^2)
\sum_{k=0}^{N-1}
q^{2k}\|v_\gamma\|^2
+
q^{2N}\|v_\gamma\|^2.
\]
El terminal\par
\[
q^{2N}\|v_\gamma\|^2
\longrightarrow0
\]
n'efface pas la hauteur : celle-ci demeure dans le caractère unitaire \(\tau_\gamma^n\).\par
Voici la lecture correcte du tour :\par
\Needspace{8\baselineskip}
\begin{center}
\small
\begin{tabularx}{\textwidth}{@{}YY@{}}
\toprule
Dato & Ce qui est enregistré \\
\midrule
\(n\) & le nombre de tours complets écoulés \\
\(q=5/9\) & la quantité de terminal restante après chaque tour \\
\(\tau_\gamma\) & le zéro individuel transporté par l'orbite \\
\(m_\gamma\) & le poids algébrique de ce zéro \\
mémoire/retenue & pourquoi le retour à une phase visible ne restitue pas le même état \\
\bottomrule
\end{tabularx}
\end{center}
C'est pourquoi il ne faut pas associer le premier zéro au premier tour. Chaque zéro traverse toute la succession des tours avec son propre caractère.\par
Une image pédagogique appropriée serait celle d'un disque qui tourne : le nombre de tours n'identifie pas la note émise. Le tour \(n\) indique le temps écoulé ; la fréquence identifie le mode. Ici, \(n\) est le compteur, \(q\) l'atténuation et \(\tau_\gamma\) la signature spectrale.\par
\Needspace{7\baselineskip}
\subsection{Localisation du champ spectral de Weil au moyen de l'entrelaceur nonadique de prolongement}
La construction précédente se situe réellement dans\par
\[
\overline{\operatorname{Ran}\mathscr L},
\]
l'espace de perte produit par la télescopie nonadique. Ce n'est pas une analogie externe ajoutée à la preuve de Riemann.\par
Il faut toutefois la distinguer de l'holonomie combinatoire native. Le TPK contient également\par
\[
M_9=qV_9,
\qquad
V_9^*V_9=I.
\]
L'identité de défaut détermine \(q\), mais ne détermine pas la partie unitaire \(V_9\). Un même bilan\par
\[
M_9^*M_9=q^2I
\]
est compatible avec de nombreuses phases unitaires différentes. La télescopie, à elle seule, ne peut indiquer laquelle contient les hauteurs.\par
L'identité entre les orbites spectrales nouvellement construites et les orbites du registre combinatoire APP–TRIT–TPK est établie par l'entrelaceur construit, sur les états cylindriques enrichis, dans le théorème de Cayley--Pontryagin de ce même chapitre :\par
\[
J:
\mathcal K_{\mathrm{loss}}
\longrightarrow
\mathcal K_{\mathrm{TPK}}^{\mathrm{unit}}
\]
tal que\par
\[
\boxed{
J\mathcal C_{\mathrm{loss}}
=
V_9J.
}
\]
En outre, \(J\) doit conserver cylindre, orientation, retenue, parcours, mémoire, frontière et multiplicité spectrale.\par
L'égalité précédente n'est ni postulée ni reçue de la réalisation classique : elle résulte du cocycle entier de mémoire, du transport cylindrique tordu par \(\mathcal C_{\mathrm{Weil}}\), de la projectivité des poids et du plongement isométrique \(J_k\) construit ci-dessous. En l'itérant, on obtient\par
\[
(qV_9)^nJv_\gamma
=
q^n\tau_\gamma^nJv_\gamma,
\]
et chaque zéro s'identifie alors au caractère unitaire d'une orbite native du TPK. L'identité conserve cylindre, orientation, retenue, parcours, mémoire, frontière et multiplicité spectrale, de sorte que la correspondance ne dépend pas d'une énumération externe zéro--tour.\par
\Needspace{7\baselineskip}
\subsection{Algorithme de calcul spectral indépendant de l'introduction préalable des zéros}
Le résultat ne se limite pas à représenter une liste déjà connue. L'opérateur \(H_W\) a été défini à partir des translations et de la forme premier–gamma–polaire. Sa résolvante peut s'écrire\par
\[
(H_W-iI)^{-1}
=
i\int_0^\infty
e^{-t}U_{-t}\,dt.
\]
Par conséquent, pour des fonctions tests \(f_j,f_k\),\par
\[
\left\langle
(H_W-iI)^{-1}[f_j],[f_k]
\right\rangle_W
=
i\int_0^\infty
e^{-t}
\mathscr I_{\mathrm{GW}}(T_{-t}f_j,f_k)\,dt.
\]
Le membre de droite se calcule à partir des canaux premier, gamma et polaire. Il ne nécessite pas l'introduction d'une table de hauteurs.\par
Puisque le spectre est discret et que \(|\gamma|\to\infty\), la résolvante est compacte. Si \(P_N\) sont des projections finies convergeant fortement vers l'identité,\par
\[
P_N(H_W-iI)^{-1}P_N
\longrightarrow
(H_W-iI)^{-1}
\]
en norma.\par
Leurs valeurs propres convergent vers\par
\[
\lambda_\gamma
=
\frac1{\gamma-i},
\]
et l'on en déduit\par
\[
\gamma=i+\lambda_\gamma^{-1}.
\]
Cela ouvre une véritable voie de calcul : former des matrices finies en utilisant uniquement la forme arithmétique construite directement à partir des données et retrouver les hauteurs comme limites spectrales.\par
Les bornes effectives de troncature et le choix de bases optimisées relèvent de la réalisation numérique de la résolvante compacte ; ils affinent son efficacité de calcul sans conditionner le théorème spectral ni l'entrelacement opératoriel déjà construit. Il n'existe pas d'énumération tautologique du type \(n=N(\gamma_n)\) : les coefficients des approximations proviennent de la forme arithmétique et leurs valeurs propres convergent vers le spectre de l'opérateur.\par
\Needspace{7\baselineskip}
\subsection{Quantification spectrale et entrelacement nonadique}
La correspondance exacte se formule au niveau typé suivant :\par
\begin{quote}
La correspondance démontrée n'identifie pas ordinalement le zéro numéro \(j\) au tour nonadique numéro \(j\), ne distribue pas neuf zéros par cycle et n'impose pas de répétition périodique des hauteurs. Elle identifie exactement le diviseur des zéros non triviaux aux atomes pondérés du générateur des translations de la forme de Weil. La transformée de Cayley convertit chaque hauteur en un caractère unitaire unique ; la réalisation nonadique de perte le transporte comme une orbite hélicoïdale \(q^n\tau_\gamma^n\) ; et le transport tordu \(\widetilde U_k\), conjointement avec \(J_k\), entrelace ce caractère spectral avec l'holonomie combinatoire native.\par
\end{quote}
L'interprétation opératorielle qui en résulte est la suivante :\par
\begin{quote}
Chaque zéro détermine le caractère angulaire individuel d'une histoire spectrale dont l'itéré traverse des tours nonadiques successifs. L'indice du tour quantifie la profondeur de mémoire et le zéro fixe le caractère unitaire qui gouverne son évolution angulaire.\par
\end{quote}
En conséquence, le résultat établit une voie opératorielle effective : il démontre la quantification spectrale individuelle des hauteurs, construit leur orbite nonadique canonique et l'identifie, au moyen de \(J_k\), au caractère de mémoire qui agit sur tous les tours du registre TPK. L'identification est harmonique et opératorielle ; elle ne dépend pas d'une numérotation des zéros par les tours.\par
\paragraph{Domaine et construction du champ spectral de Weil}
La forme positive \(\mathscr I_{\mathrm{GW}}=\mathscr L^*\mathscr L\), la contraction \(q=5/9\), la télescopie et le tour \(M_9=qV_9\) constituent la structure mathématique initiale. Sur celle-ci se construisent l'espace de Hilbert de Weil, le groupe des translations, son générateur de Stone, la correspondance de Cayley et les orbites \(q^n\tau_\gamma^n\). L'approximation de la résolvante directement à partir de la forme arithmétique définit sa réalisation informatique. Le cocycle de mémoire et le transport tordu construits ci-après fournissent l'entrelaceur exact avec l'holonomie combinatoire native.

\section{Théorème de Cayley–Pontryagin du double cercle}
Il n'existe pas de correspondance « zéro numéro \(j\) = tour nonadique numéro \(j\) ». Ce qui existe est plus profond : chaque zéro détermine un caractère unitaire de la mémoire verticale du registre, et l'orbite complète de ce caractère est exactement l'orbite spectrale du zéro après la transformation de Cayley.\par
\Needspace{7\baselineskip}
\subsection{Théorème du double cercle}
La connexion nonadique possède une mémoire verticale entière\par
\[
q_{\mathrm{mem}}:\operatorname{Obj}(\mathrm{TPK})\longrightarrow\mathbb Z.
\]
Pour tout chemin admissible \(\gamma:x\to y\), on définit le cocycle\par
\[
m(\gamma)
=
q_{\mathrm{mem}}(y)-q_{\mathrm{mem}}(x).
\]
Es aditivo:\par
\[
m(\gamma_2\gamma_1)
=
m(\gamma_2)+m(\gamma_1),
\]
et pour tout tour complet de la connexion,\par
\[
m(\Gamma_9)=1.
\]
Par itération,\par
\[
m(\Gamma_9^n)=n.
\]
C'est la donnée décisive. La phase visible revient après neuf pas, mais l'état ne revient pas au même point : la mémoire a avancé d'une unité.\par
D'autre part, la démonstration de Riemann construit d'abord la forme positive de Weil et son espace GNS :\par
\[
\mathcal H_W
=
\overline{\mathcal D/\ker \mathscr I_{\mathrm{GW}}}.
\]
La translation construite à partir de la forme arithmétique produit un générateur autoadjoint \(H_W\). Sans introduire les zéros comme données, on forme sa transformée de Cayley :\par
\[
\mathcal C_{\mathrm{Weil}}
=
\left(H_W+\frac{i}{2}I\right)
\left(H_W-\frac{i}{2}I\right)^{-1}.
\]
Puisque \(H_W\) est autoadjoint,\par
\[
\mathcal C_{\mathrm{Weil}}^*
\mathcal C_{\mathrm{Weil}}=I.
\]
Par conséquent, \(\mathcal C_{\mathrm{Weil}}\) est unitaire et son spectre appartient à \(S^1\).\par
Après la reconnaissance spectrale, et non avant, une hauteur\par
\[
\rho_\gamma=\frac12+i\gamma
\]
produit la phase\par
\[
\tau_\gamma
=
\frac{\gamma+i/2}{\gamma-i/2}
=
\frac{\rho_\gamma-1}{\rho_\gamma},
\qquad
|\tau_\gamma|=1.
\]
L'inversion est exacte :\par
\[
\rho_\gamma=\frac{1}{1-\tau_\gamma},
\qquad
\gamma
=
\frac12\cot\left(\frac{\arg\tau_\gamma}{2}\right).
\]
Ainsi apparaissent deux réalisations circulaires obtenues par des lecteurs distincts du même état enrichi :\par
\[
\left\{\Re s=\frac12\right\}\cup\{\infty\}
\overset{\tau(s)=(s-1)/s}{\cong}
S^1,
\]
y\par
\[
\widehat{\mathbb Z_{\mathrm{mem}}}
\cong S^1.
\]
Le premier est le cercle de Cayley de la droite critique. Le second est le dual de Pontryagin de la mémoire verticale du registre. L'entrelaceur démontre qu'il ne s'agit pas de deux cercles simplement semblables : ce sont les deux réalisations du même opérateur de phase.\par
\Needspace{7\baselineskip}
\subsection{Isométrie de transport tordu \(\widetilde U_k\) sur les états cylindriques, les poids projectifs et la mémoire de Weil}
Soit \(Z_k\) l'ensemble des états cylindriques enrichis à la profondeur \(k\), muni d'une mesure projective de support plein \(\mu_k\), et soit\par
\[
\tau_{k+1,k}:Z_{k+1}\longrightarrow Z_k
\]
el truncamiento.\par
Pour un descendant \(z'\) de \(z\), on définit\par
\[
w_k(z'|z)
=
\sqrt{\frac{\mu_{k+1}(z')}{\mu_k(z)}}.
\]
La projectivité donne\par
\[
\sum_{\tau_{k+1,k}z'=z}w_k(z'|z)^2=1.
\]
Sur\par
\[
\mathscr K_k
=
\ell^2(Z_k)\otimes\mathcal H_W
\]
on définit le transport tordu\par
\[
\widetilde U_k(e_z\otimes v)
=
\sum_{\tau_{k+1,k}z'=z}
w_k(z'|z)\,
e_{z'}\otimes
\mathcal C_{\mathrm{Weil}}^{\,m(z\to z')}v.
\]
Sous forme matricielle,\par
\[
[\widetilde U_k]_{z',z}
=
\mathbf 1_{\{\tau z'=z\}}
\sqrt{\frac{\mu_{k+1}(z')}{\mu_k(z)}}
\,
\mathcal C_{\mathrm{Weil}}^{\,m(z\to z')}.
\]
Cette formule conserve matériellement tous les descendants. Elle n'additionne pas des chemins distincts comme s'ils étaient identiques : registre, orientation, retenue, frontière, mémoire, parcours et terminal restent présents dans l'étiquette \(z'\).\par
L'isométrie se démontre directement. Les descendants de prédécesseurs différents sont orthogonaux ; pour chaque prédécesseur, la somme des carrés des poids vaut un ; et chaque puissance de \(\mathcal C_{\mathrm{Weil}}\) est unitaire. Par conséquent,\par
\[
\widetilde U_k^*\widetilde U_k=I.
\]
La composition est également exacte :\par
\[
\widetilde U_{\ell,j}\widetilde U_{j,k}
=
\widetilde U_{\ell,k},
\]
car les poids se télescopent et le cocycle de mémoire est additif.\par
Il existe même une équivalence de jauge explicite. Si\par
\[
G_k(e_z\otimes v)
=
e_z\otimes
\mathcal C_{\mathrm{Weil}}^{\,q_{\mathrm{mem}}(z)}v,
\]
alors\par
\[
\widetilde U_k
=
G_{k+1}(U_k\otimes I)G_k^*.
\]
Par conséquent, le système tordu n'est pas un ornement arbitraire : c'est le transport cylindrique natif vu dans la jauge spectrale déterminée par la mémoire.\par
Définissons maintenant le vecteur de mesure\par
\[
\Omega_k
=
\sum_{z\in Z_k}
\sqrt{\mu_k(z)}\,e_z,
\qquad
\|\Omega_k\|=1,
\]
et le plongement\par
\[
J_k:\mathcal H_W\longrightarrow\mathscr K_k,
\qquad
J_kv=\Omega_k\otimes v.
\]
Au cours d'un tour complet, tous les parcours présentent un accroissement de mémoire égal à un. Ainsi,\par
\[
\boxed{
\widetilde U_{\Gamma_9,k}J_k
=
J_{k+9}\mathcal C_{\mathrm{Weil}}
}
\]
et, après \(n\) tours,\par
\[
\boxed{
\widetilde U_{\Gamma_9^n,k}J_k
=
J_{k+9n}\mathcal C_{\mathrm{Weil}}^n.
}
\]
Tel est l'entrelaceur annoncé et construit au sein de la même chaîne causale.\par
\Needspace{7\baselineskip}
\subsection{Identité orbitale sous la conjugaison entre les \(2\) réalisations spectrales}
Soit \(\psi_\gamma\) un vecteur de la fibre spectrale associée à \(\gamma\). Alors\par
\[
\mathcal C_{\mathrm{Weil}}\psi_\gamma
=
\tau_\gamma\psi_\gamma.
\]
En appliquant le carré précédent :\par
\[
\widetilde U_{\Gamma_9^n}J\psi_\gamma
=
J\mathcal C_{\mathrm{Weil}}^n\psi_\gamma
=
\tau_\gamma^nJ\psi_\gamma.
\]
Or le caractère de mémoire correspondant à \(\tau_\gamma\) est\par
\[
\chi_{\tau_\gamma}(m)
=
\tau_\gamma^m.
\]
Comme\par
\[
m(\Gamma_9^n)=n,
\]
se obtiene\par
\[
\boxed{
\widetilde U_{\Gamma_9^n}J\psi_\gamma
=
\chi_{\tau_\gamma}
\!\left(m(\Gamma_9^n)\right)
J\psi_\gamma.
}
\]
C'est l'égalité littérale recherchée.\par
Cela ne signifie pas que le zéro numéro 37 soit le tour numéro 37. Cela signifie que chaque zéro sélectionne un caractère du groupe complet de mémoire et que ce caractère parcourt tous les tours :\par
\[
n\longmapsto\tau_\gamma^n.
\]
Por tanto:\par
\begin{itemize}
\item le tour nonadique fournit la variable entière \(n\) ;
\item la mémoire fournit le groupe \(\mathbb Z\) ;
\item le zéro fournit un caractère \(\chi_{\tau_\gamma}\in\widehat{\mathbb Z}\) ;
\item l'orbite spectrale et l'orbite du registre coïncident terme à terme après le plongement \(J\).
\end{itemize}
La réserve précédente est levée sans affirmer « neuf zéros par cycle ». Il n'y a pas d'énumération zéro–tour. Il existe une décomposition harmonique de la même dynamique de tour.\par
\Needspace{7\baselineskip}
\subsection{Compactification circulaire et relèvement spiral de l'orbite spectrale}
Une distinction cruciale doit être conservée. La phase est unitaire :\par
\[
\mathcal C_{\mathrm{Weil}}^*
\mathcal C_{\mathrm{Weil}}=I.
\]
La dynamique stricte du feuillet de perte est\par
\[
M_9=qV_9,
\qquad
q=\frac59.
\]
L'opérateur conjoint correct est\par
\[
\widetilde M_9
=
\frac59\,\widetilde U_{\Gamma_9}.
\]
Sur la fibre spectrale,\par
\[
\widetilde M_9^nJ\psi_\gamma
=
\left(\frac59\right)^n
\tau_\gamma^n
J\psi_\gamma.
\]
Ainsi apparaissent deux images simultanées :\par
\begin{itemize}
\item à la frontière normalisée, l'orbite est circulaire :
\end{itemize}
  \[
  \tau_\gamma^n\in S^1;
  \]
\begin{itemize}
\item dans la réalisation stricte, elle est une spirale :
\end{itemize}
  \[
  \left(\frac59\right)^n\tau_\gamma^n\longrightarrow0.
  \]
Le zéro terminal n'efface ni la phase ni la généalogie. Le rayon s'éteint, mais le registre conserve le nombre de tours effectués et le caractère transporté.\par
La télescopie est exacte :\par
\[
I
=
\sum_{r=0}^{L-1}
\widetilde M_9^{*r}
D_9^*D_9
\widetilde M_9^r
+
\widetilde M_9^{*L}\widetilde M_9^L,
\]
con\par
\[
D_9^*D_9
=
I-\widetilde M_9^*\widetilde M_9
=
\frac{56}{81}I,
\]
y terminal\par
\[
\widetilde M_9^{*L}\widetilde M_9^L
=
\left(\frac{25}{81}\right)^LI
\longrightarrow0.
\]
Cela répond par avance à une objection de typage : une contraction n'a pas été confondue avec un opérateur unitaire. La phase appartient à \(\mathcal C_{\mathrm{Weil}}\) ; la perte radiale appartient à \(5/9\).\par
\Needspace{7\baselineskip}
\subsection{Compatibilité entre cylindres généalogiques, génération des constantes et entrelaceur spectral}
Ton intuition concernant les cylindres est correcte, sous réserve d'une distinction précise.\par
Les cylindres numériques sont semi-ouverts :\par
\[
C(N_K,6K)
=
\left[
\frac{N_K}{3^{6K}},
\frac{N_K+1}{3^{6K}}
\right),
\]
et chacun se divise en 729 sous-cylindres. Leur diamètre satisfait\par
\[
\operatorname{diam}C(N_K,6K)=3^{-6K}\longrightarrow0.
\]
Les cylindres généalogiques \([a]_k\) de la limite inverse sont des boules ouvertes et fermées de l'ultramétrique :\par
\[
d_\vartheta(x,y)
=
\vartheta^{N(x,y)}.
\]
Par conséquent, « rayon tendant vers zéro » possède deux expressions :\par
\begin{enumerate}
\item Dans la lecture archimédienne, l'intersection d'intervalles emboîtés exprime un réel unique.
\item Dans l'état enrichi, l'histoire complète continue de conserver tous les préfixes qui ont produit ce réel.
\end{enumerate}
Ainsi, conformément à la \hyperref[sec:tesis-inaugural-helicoidal-hmt-md]{tesis inaugural}, le même système coinductif exprime de manière corrélée \(\pi\), \(e\), \(\varphi\) et \(\alpha\). Dans l'ordre opératoriel interne, le registre \(U_{12}^{\mathrm{sgn}}\), le sceau \(K\) et la retenue dodécaphasique interviennent une fois \(P,E,\Phi\) disponibles : ils expriment et certifient la coordonnée \(\alpha\) de cette même génération coinductive corrélée, sans l'incorporer comme objet causalement externe. Les sections archimédiennes et l'expression signée conservent des codomaines différents sans rompre leur généalogie commune.\par
Ces constantes ne sélectionnent pas les phases de Riemann. Leur importance est plus profonde : elles démontrent que le même système de cylindres finis peut produire un objet continu sans détruire l'histoire qui l'a engendré.\par
Le nouvel entrelaceur utilise exactement cette propriété. Le cercle spectral ne se superpose pas au continu depuis l'extérieur : il apparaît comme le dual de sa mémoire entière.\par
\Needspace{7\baselineskip}
\subsection{Invariant gradué du corridor spectral \(4\to2\to6\to54\)}
La chaîne \(4\to54\) peut maintenant être située sans numérologie.\par
Le bloc central d'APP est\par
\[
B_c
=
\begin{pmatrix}
7&2\\
2&7
\end{pmatrix}
=
9P_++5P_-.
\]
Ses valeurs propres sont \(9\) et \(5\), de sorte que\par
\[
9-5=4,
\qquad
q=\frac59.
\]
La hiérarchie typée est\par
\[
\underbrace{9-5}_{4};
\qquad
\underbrace{2}_{\text{couplage}};
\qquad
\underbrace{51-45}_{6};
\qquad
\underbrace{9\cdot6}_{54}.
\]
Dans la connexion,\par
\[
\Gamma_{54}=\Gamma_9^6.
\]
Par le nouvel entrelaceur,\par
\[
\widetilde U_{\Gamma_{54}}J
=
J\mathcal C_{\mathrm{Weil}}^6,
\]
et, dans la fibre \(\gamma\),\par
\[
\widetilde U_{\Gamma_{54}}J\psi_\gamma
=
\tau_\gamma^6J\psi_\gamma.
\]
Dans la réalisation stricte,\par
\[
\widetilde M_{54}J\psi_\gamma
=
\left(\frac59\right)^6
\tau_\gamma^6J\psi_\gamma.
\]
C'est l'incidence spectrale exacte de la profondeur 54 : six tours nonadiques, six unités de mémoire et sixième puissance du caractère.\par
Elle doit être distinguée d'autres objets dont la valeur est également 54 :\par
\begin{itemize}
\item le défaut résiduel APP ;
\item la trace orientée APP–Fock ;
\item la norme du vecteur radial de \(A_2^{12}\) ;
\item le coefficient du caractère \(t_3\) ;
\item la valeur
\end{itemize}
  \[
  54=\operatorname{Tr}(3B\mid V^\natural_2).
  \]
Le fait qu'ils valent tous 54 possède une importance structurelle, mais ne les transforme pas automatiquement en un même morphisme.\par
\Needspace{7\baselineskip}
\subsection{Leech, \(V^\natural\) et le Monstre}
La chaîne exceptionnelle demeure exacte :\par
\[
A_2^{12}
\longrightarrow
C_W=[12,6,6]_3
\longrightarrow
N_\alpha\cong\Lambda_{24}
\longrightarrow
V^\natural
\longrightarrow
\mathbb M.
\]
La rigidité transversale sélectionne\par
\[
m=12,
\qquad
F(12)=R(12)=54.
\]
L'isométrie sans point fixe d'ordre trois produit l'orbifold\par
\[
\operatorname{Orb}_{C_3}
(V_{N_\alpha},\widehat g_c)
\cong V^\natural,
\]
et l'automorphisme dual appartient à la classe \(3B\) :\par
\[
T_{3B}
=
q^{-1}+54q+\cdots.
\]
Cependant, le Monstre n'engendre pas les phases \(\tau_\gamma\). Il est fini, tandis que les caractères de mémoire parcourent tout \(S^1\). En particulier, un élément \(3B\) est d'ordre trois, mais une phase spectrale générale n'est pas une racine cubique de l'unité.\par
La réunion correcte est un produit de représentations :\par
\[
\mathscr K^{\mathrm{gold}}
=
\mathscr H_{\mathrm{cyl}}
\otimes
\mathcal H_W
\otimes
\mathcal H_{V^\natural}.
\]
La mémoire et la transformation de Cayley agissent sur les deux premiers facteurs ; le Monstre agit sur le troisième :\par
\[
V_{\mathrm{mem}}
=
\widetilde U_{\Gamma_9}\otimes I,
\qquad
\rho_{\mathbb M}(g)
=
I\otimes I\otimes\rho_{V^\natural}(g).
\]
Les deux actions commutent. Ainsi, chaque orbite spectrale peut conserver simultanément :\par
\begin{itemize}
\item su altura \(\gamma\);
\item sa phase \(\tau_\gamma\) ;
\item sa mémoire nonadique \(n\) ;
\item sa multiplicité spectrale ;
\item son type transversal du Monstre.
\end{itemize}
Mais la multiplicité d'un zéro ne s'identifie pas par décret à une dimension monstrueuse. Cette identification exigerait un entrelaceur supplémentaire entre la fibre spectrale et une représentation concrète du Monstre.\par
Une véritable classe calculable de traces entrelacées demeure en revanche ouverte :\par
\[
\Theta_g(n;a,\beta)
=
\operatorname{Tr}
\left(
e^{-aH_W^2}\mathcal C_{\mathrm{Weil}}^n
\right)
\operatorname{Tr}_{V^\natural}
\left(
g\,e^{-\beta(L_0-1)}
\right).
\]
Après la reconnaissance spectrale,\par
\[
\Theta_g(n;a,\beta)
=
\left(
\sum_\gamma
m_\gamma e^{-a\gamma^2}\tau_\gamma^n
\right)
T_g(e^{-\beta}).
\]
C'est bien une nouvelle voie RH–Moonshine rigoureusement typée : moments des orbites des zéros, mémoire nonadique et caractères monstrueux au sein d'une même observable régularisée.\par
\Needspace{7\baselineskip}
\subsection{Approximation finie du champ spectral par troncatures matricielles}
La construction possède une matrice finie naturelle.\par
Pour tout opérateur unitaire \(C\), définissons la couture à neuf phases\par
\[
\mathbb U_{C,9}
=
\sum_{r=1}^{8}
|r+1\rangle\langle r|\otimes I
+
|1\rangle\langle9|\otimes C.
\]
Alors\par
\[
\mathbb U_{C,9}^9
=
I\otimes C.
\]
Con\par
\[
u_9=\frac1{3}(1,\ldots,1),
\qquad
\iota_9v=u_9\otimes v,
\]
la compression visible est\par
\[
A_{C,9}
=
\iota_9^*\mathbb U_{C,9}\iota_9
=
\frac{8I+C}{9},
\]
et le défaut satisfait\par
\[
c_{C,9}^*c_{C,9}
=
\frac{8}{81}
(I-C)^*(I-C).
\]
En prenant \(C=\mathcal C_{\mathrm{Weil}}\), on obtient une matrice à neuf phases construite sans introduire de zéros.\par
Pour rendre également finie la fibre de Weil, on choisit des fonctions tests \(f_1,\ldots,f_N\) et l'on calcule, à partir des canaux premier–gamma–polaire, les matrices\par
\[
G_N
=
\bigl(
\mathscr I_{\mathrm{GW}}(f_i,f_j)
\bigr)_{i,j},
\]
y\par
\[
A_N
=
\bigl(
\langle H_Wf_j,f_i\rangle_W
\bigr)_{i,j}.
\]
Après quotient par \(\ker G_N\), le problème généralisé\par
\[
A_Nv=\gamma_NG_Nv
\]
produit l'approximation spectrale, et sa transformée de Cayley finie est\par
\[
C_N
=
\left(
G_N^{-1}A_N+\frac{i}{2}I
\right)
\left(
G_N^{-1}A_N-\frac{i}{2}I
\right)^{-1}.
\]
Ses valeurs propres \(\tau_N\in S^1\) permettent de retrouver\par
\[
\gamma_N
=
\frac12
\cot\left(
\frac{\arg\tau_N}{2}
\right).
\]
La matrice totale des cylindres est alors\par
\[
[\widetilde U_{k,N}]_{z',z}
=
\mathbf1_{\{\tau z'=z\}}
\sqrt{\frac{\mu_{k+1}(z')}{\mu_k(z)}}
\,C_N^{\,m(z\to z')}.
\]
Cela transforme la voie conceptuelle en un programme numérique concret :\par
\begin{enumerate}
\item construire \(G_N\) et \(A_N\) sans introduire de zéros ;
\item former \(C_N\) ;
\item l'insérer dans la couture nonadique ;
\item diagonaliser les phases ;
\item retrouver les hauteurs ;
\item vérifier la stabilité sous l'augmentation simultanée de \(N\) et de la profondeur cylindrique.
\end{enumerate}
La convergence rigoureuse exige de démontrer que les sous-espaces finis forment un cœur de \(H_W\) et que les approximations convergent au sens de la résolvante forte. Il s'agit d'un certificat numérique ultérieur ; il n'affecte pas l'entrelaceur infini, qui est déjà démontré.\par
\Needspace{7\baselineskip}
\subsection{Indépendance de l'entrelaceur spectral vis-à-vis des zéros, des contractions, des multiplicités de parcours et des symétries transversales}
\textbf{« Les zéros ont été introduits pour fabriquer l'application. »} Non. \(\mathcal H_W\), \(H_W\), sa résolvante et \(\mathcal C_{\mathrm{Weil}}\) se construisent à partir de la forme premier–gamma–polaire. Les zéros apparaissent ensuite pour reconnaître le spectre.\par
\textbf{« Une contraction a été identifiée à un opérateur unitaire. »} Non. La phase est \(\mathcal C_{\mathrm{Weil}}\) ; la contraction est \((5/9)\mathcal C_{\mathrm{Weil}}\).\par
\textbf{« Des chemins distincts ont été effacés. »} Non. Chaque descendant conserve son étiquette complète dans \(e_{z'}\) ; l'orthogonalité du registre empêche de confondre deux parcours.\par
\textbf{« On affirme qu'il existe neuf zéros par cycle. »} Non. Chaque zéro détermine un caractère évalué sur tous les tours.\par
\textbf{« Le Monstre produit les hauteurs. »} Non. Le cercle des phases provient de \(\widehat{\mathbb Z_{\mathrm{mem}}}\). Le Monstre organise la symétrie transversale et commute avec cette action.\par
\textbf{« Les différents 54 constituent le même objet. »} Non. Défaut, profondeur, norme et valeur du caractère restent distincts. Leurs coïncidences sont reliées par des applications, non par une simple égalité numérique.\par
\textbf{« On a seulement renommé un opérateur préalable \(V_9\). »} Non. Une réalisation explicite propre au codomaine de \(\Gamma_9\) a été construite. Pour l'identifier à un représentant ambiant préalable quelconque de \(V_9\), fixé seulement à équivalence unitaire près, la jauge antérieure doit être choisie de façon compatible. Dans le sous-espace cyclique pertinent, la classe unitaire est déjà déterminée.\par
\Needspace{7\baselineskip}
\subsection{Domaine de la conjugaison Cayley–Pontryagin et théorème du double cercle}
APP–TRIT–TPK, les cylindres, la mesure projective, le cocycle de mémoire, \(\Gamma_9\), la contraction \(5/9\), la chaîne \(4\to54\), les constantes et le corridor exceptionnel constituent la structure initiale. Sur ce domaine se construit le transport cylindrique tordu par \(\mathcal C_{\mathrm{Weil}}^{m(\gamma)}\). Le théorème du double cercle établit l'entrelacement exact
  \[
  \widetilde U_{\Gamma_9}J
  =
  J\mathcal C_{\mathrm{Weil}}.
  \]
La certification de convergence des matrices finies et l'étude non factorisée des traces RH–Moonshine constituent la frontière ouverte après le théorème.
Par conséquent, la conclusion forte est :\par
\[
\boxed{
\text{les zéros ne sont pas les tours ;
ils sont les caractères spectraux de la mémoire des tours.}
}
\]
Et, après le plongement canonique,\par
\[
\boxed{
\text{l'orbite spectrale du zéro et l'orbite du registre sont littéralement la même orbite.}
}
\]

\section{Reconstruction effective du champ spectral de Weil}
\Needspace{5\baselineskip}
\subsection{Construction du champ spectral indépendante de données préalables sur les zéros}
On adopte la classe de fonctions à support compact\par
\[
\beta_A(x)=Z_A^{-1}
\exp\!\left[-\frac{1}{1-(x/A)^2}\right]
\mathbf 1_{\{|x|<A\}},
\qquad
f_{\omega,A}(x)=\beta_A(x)e^{-i\omega x}.
\]
La modulation \(\omega\) n'a pas été choisie à l'aide des zéros. On a préalablement balayé l'intervalle\par
\[
5\leq\omega\leq50,
\qquad
\Delta\omega=0.025.
\]
Pour chaque \(\omega\), on calcule\par
\[
M_A(\omega)
=
\mathscr I_{\rm GW}
(f_{\omega,A},f_{\omega,A})
\]
au moyen de la formule construite à partir des données arithmétiques\par
\[
\begin{aligned}
\mathscr I_{\rm GW}(f,g)
={}&c_\Gamma C_{f,g}(0)
+\int_0^\infty
\mu_\Gamma(a)
\bigl[
2C_{f,g}(0)-C_{f,g}(a)-C_{f,g}(-a)
\bigr]\,da\\
&-\sum_{p,k}
\frac{\log p}{p^{k/2}}
\bigl[
C_{f,g}(k\log p)+C_{f,g}(-k\log p)
\bigr]\\
&+P_f^+\overline{P_g^-}
+P_f^-\overline{P_g^+},
\end{aligned}
\]
où\par
\[
C_{f,g}(a)=\int_{\mathbb R}f(x)\overline{g(x-a)}\,dx,
\]
\[
\mu_\Gamma(a)=\frac{e^{-a/2}}{1-e^{-2a}},
\qquad
c_\Gamma=-\gamma_{\!E}-\frac{\pi}{2}-3\log2-\log\pi,
\]
y\par
\[
P_f^\pm=\int_{\mathbb R}f(x)e^{\pm x/2}\,dx.
\]
Ni \(\zeta(s)\), ni une routine de calcul des zéros, ni Riemann–Siegel, ni Odlyzko, ni une table de hauteurs ne sont intervenus.\par
Avec \(A=3\), le premier balayage à l'aveugle a produit dix maxima dominants approximativement en\par
\[
14.125,\ 21.025,\ 25.000,\ 30.425,\ 32.925,\ 37.575,\ 40.925,\ 43.325,\ 48.000,\ 49.775.
\]
Les lobes latéraux avaient une masse environ cent fois plus faible.\par
\Needspace{5\baselineskip}
\subsection{Raffinement du champ spectral indépendant des valeurs cibles}
J'ai augmenté la largeur compacte \(A\). Cela augmente simultanément la résolution spectrale et le nombre de puissances de nombres premiers utilisées :\par
\Needspace{8\baselineskip}
\begin{center}
\small
\begin{tabularx}{\textwidth}{@{}RRRRR@{}}
\toprule
\(n\) & \(A=3\) & \(A=4\) & \(A=5\) & \(A=6\) \\
\midrule
1 & 14.134699719 & 14.134721825 & 14.134724556 & 14.134725023 \\
2 & 21.022713497 & 21.022092670 & 21.022033480 & 21.022034382 \\
3 & 25.010192460 & 25.010816662 & 25.010862963 & 25.010863775 \\
4 & 30.431966224 & 30.424255364 & 30.424903280 & 30.424992486 \\
5 & 32.927889053 & 32.935726409 & 32.935037177 & 32.934942801 \\
6 & 37.585309844 & 37.585911673 & 37.586134491 & 37.586172107 \\
7 & 40.926249022 & 40.917939546 & 40.919104383 & 40.918785046 \\
8 & 43.320460106 & 43.328124621 & 43.326729729 & 43.327013819 \\
9 & 47.996518745 & 48.009530251 & 48.005271807 & 48.005312108 \\
10 & 49.781332020 & 49.769281661 & 49.773747976 & 49.773697939 \\
\bottomrule
\end{tabularx}
\end{center}
Dans \(A=6\) sont entrées exactement 15.040 puissances de nombres premiers \(p^k\), avec \(p^k\leq e^{12}\). Pour le premier maximum, le registre numérique était :\par
\[
\begin{array}{lr}
\text{canal premier} & 8.077017879363272,\\
\text{intégrale gamma} & 6.182517531024263,\\
c_\Gamma I & -5.372183419225665,\\
\text{canal polaire} & -1.0505899\times10^{-8},\\ \hline
\text{masse totale} & 8.887351980655973.
\end{array}
\]
La queue gamma au-delà de \(a=65\) a été bornée par\par
\[
1.54\times10^{-14}.
\]
Cela montre un fait conceptuel : le zéro n'est pas introduit dans le calcul. Il apparaît comme le centre stable où les quatre composantes de la fonctionnelle se referment.\par
\Needspace{5\baselineskip}
\subsection{Approximations matricielles finies de l'opérateur spectral}
Une matrice de Gram isolée ne contient pas de hauteurs identifiables. Ses valeurs propres dépendent de la base et représentent des masses ou des énergies, non des positions spectrales.\par
C'est pourquoi on construit trois matrices :\par
\[
G_{ij}=\mathscr I_{\rm GW}(f_j,f_i),
\]
\[
K_{ij}=\mathscr I_{\rm GW}(H_Wf_j,f_i),
\]
\[
Q_{ij}=\mathscr I_{\rm GW}(H_Wf_j,H_Wf_i).
\]
Avec l'orientation positive utilisée dans le calcul,\par
\[
H_W=i\,\partial_x
\]
sur la classe réfléchie \(f_{\omega,A}=\beta_Ae^{-i\omega x}\). L'orientation conjuguée \(-i\partial_x\) restitue le feuillet spectral opposé ; les deux produisent les mêmes hauteurs positives par la symétrie \(\gamma\leftrightarrow-\gamma\).\par
Ce générateur n'est pas introduit arbitrairement. La positivité démontrée permet de former la construction GNS de \(\mathscr I_{\rm GW}\). Les translations sont des isométries fortement continues de cette construction GNS, et Stone produit leur générateur autoadjoint. Sur le cœur des fonctions tests, la dérivation des translations est précisément \(i\partial_x\).\par
Les trois matrices ont été calculées au moyen du même évaluateur premier–gamma–polaire. Ni \(K\) ni \(Q\) n'ont été construits à partir d'une liste spectrale.\par
Dans le bloc \(10\times10\), les défauts d'hermiticité avant symétrisation numérique étaient\par
\[
\|G-G^*\|_2=2.11\times10^{-14},
\]
\[
\|K-K^*\|_2=7.14\times10^{-13},
\]
\[
\|Q-Q^*\|_2=1.84\times10^{-10}.
\]
Les valeurs propres de \(G\) étaient\par
\[
\begin{aligned}
4.90806959,\ 5.50537820,\ 5.61057540,\ 5.88666432,\ 5.91920889,\\
5.93049827,\ 5.96524271,\ 6.27603538,\ 6.33591428,\ 7.02999383.
\end{aligned}
\]
Elles sont toutes positives et le conditionnement du bloc n'est que\par
\[
\kappa(G)=1.43233.
\]
Il n'a pas été nécessaire de dissimuler une valeur propre négative ni de projeter artificiellement une partie gênante.\par
La métrique a été blanchie :\par
\[
\widehat H_N
=
G_N^{-1/2}K_NG_N^{-1/2},
\]
et l'on a résolu le faisceau hermitien\par
\[
K_Nv=\gamma_NG_Nv.
\]
C'est ce qui transforme les masses de Gram en hauteurs.\par
\Needspace{5\baselineskip}
\subsection{Transformée de Cayley et compactification circulaire initiale du spectre}
À chaque \(\widehat H_N\), on a appliqué\par
\[
C_N=
\left(\widehat H_N+\frac{i}{2}I\right)
\left(\widehat H_N-\frac{i}{2}I\right)^{-1}.
\]
Si \(\gamma\) est une valeur propre du générateur, sa phase est\par
\[
\tau_\gamma
=
\frac{\gamma+i/2}{\gamma-i/2},
\qquad |\tau_\gamma|=1.
\]
La reconstruction inverse est\par
\[
\gamma
=
\frac{i}{2}\frac{\tau+1}{\tau-1}
=
\frac12\cot\!\left(\frac{\arg\tau}{2}\right).
\]
Le défaut d'unitarité calculé était\par
\[
\|C_N^*C_N-I\|_2=1.45\times10^{-15}.
\]
Les trois premières phases étaient\par
\[
\tau_1=
0.997500505583064
+0.070659333152323\,i,
\]
\[
\tau_2=
0.998869228927548
+0.047542228615055\,i,
\]
\[
\tau_3=
0.999201013749813
+0.039966662624574\,i.
\]
C'est le premier cercle : la droite spectrale est compactifiée sans perdre aucune hauteur, car la transformation de Cayley est bijective entre \(\mathbb R\) et le cercle privé du point terminal.\par
\Needspace{5\baselineskip}
\subsection{Immersion exacte du spectre dans le cercle nonadique}
Une fois \(C_N\) construit, on définit la matrice de couture à neuf positions\par
\[
S_{N,9}
=
\sum_{r=1}^{8}
|r+1\rangle\langle r|\otimes I
+
|1\rangle\langle9|\otimes C_N.
\]
On n'associe pas un zéro à chaque porte. Les neuf positions constituent le calendrier de transport. La phase spectrale ne s'applique que lorsque le parcours achève le tour.\par
L'identité exacte est\par
\[
S_{N,9}^{\,9}=I_9\otimes C_N.
\]
Par conséquent, si \(J_Nv=u_9\otimes v\), l'holonomie du tour complet satisfait\par
\[
S_{N,9}^{\,9}J_N=J_NC_N.
\]
C'est l'entrelacement du double cercle :\par
\begin{itemize}
\item le cercle de Cayley contient la phase spectrale ;
\item le cercle nonadique transporte l'état pendant neuf pas ;
\item un tour complet applique exactement une fois la phase de Cayley ;
\item la mémoire enregistre le nombre de tours ;
\item aucun zéro ne s'identifie à une porte isolée.
\end{itemize}
Pour les dix candidats, une matrice \(90\times90\) a été construite. Les contrôles étaient :\par
\[
\|S_{N,9}^*S_{N,9}-I\|_2
=
4.44\times10^{-16},
\]
\[
\|S_{N,9}^{\,9}-I_9\otimes C_N\|_2=0.
\]
La contraction visible est\par
\[
M_N=\frac59S_{N,9},
\qquad
D_N=\frac{\sqrt{56}}9I.
\]
On a vérifié\par
\[
M_N^*M_N=\frac{25}{81}I,
\qquad
D_N^*D_N=\frac{56}{81}I,
\]
ainsi que la télescopie\par
\[
I=
\sum_{r=0}^{L-1}
M_N^{*r}D_N^*D_NM_N^r
+
M_N^{*L}M_N^L.
\]
Para \(L=20\),\par
\[
\left\|
\sum_{r=0}^{19}M_N^{*r}D_N^*D_NM_N^r
+
M_N^{*20}M_N^{20}
-I
\right\|_2
=
3.33\times10^{-16},
\]
tandis que le terminal était\par
\[
\|M_N^{*20}M_N^{20}\|_2
=
6.1531824951\times10^{-11},
\]
exactement égal, à la précision machine près, à\par
\[
\left(\frac{25}{81}\right)^{20}.
\]
La perte visible ne détruit pas la phase : celle-ci appartient à la dilatation unitaire, et le facteur \(5/9\) règle la quantité qui reste visible à chaque tour.\par
\Needspace{5\baselineskip}
\subsection{Reconstruction des hauteurs et des résidus avec comparaison externe ultérieure}
Pour chaque vecteur généralisé \(v\), le résidu a été calculé sans zéros au moyen de\par
\[
r_N(\gamma,v)^2
=
\frac{
v^*
\left(
Q_N-2\gamma K_N+\gamma^2G_N
\right)v
}{
v^*G_Nv
}.
\]
Ce résidu mesure l'écart entre le vecteur fini et un véritable vecteur spectral. Une hauteur n'est pas acceptée au seul motif que le faisceau produit un nombre.\par
\Needspace{8\baselineskip}
\begin{center}
\small
\begin{tabularx}{\textwidth}{@{}RRRR@{}}
\toprule
\(n\) & altura producida & résidu & erreur dans la comparaison ultérieure \\
\midrule
1 & 14.134725141525 & \(8.15\times10^{-5}\) & \(2.10\times10^{-10}\) \\
2 & 21.022039638825 & \(4.82\times10^{-5}\) & \(5.34\times10^{-11}\) \\
3 & 25.010857580595 & \(1.15\times10^{-4}\) & \(4.49\times10^{-10}\) \\
4 & 30.424876128079 & \(2.32\times10^{-4}\) & \(2.22\times10^{-9}\) \\
5 & 32.935061595211 & \(3.98\times10^{-4}\) & \(7.47\times10^{-9}\) \\
6 & 37.586178242423 & \(1.18\times10^{-3}\) & \(8.36\times10^{-8}\) \\
7 & 40.918719778888 & \(3.14\times10^{-3}\) & \(7.67\times10^{-7}\) \\
8 & 43.327078423026 & \(7.18\times10^{-3}\) & \(5.14\times10^{-6}\) \\
9 & 48.005155742413 & \(6.22\times10^{-3}\) & \(4.86\times10^{-6}\) \\
10 & 49.773888321513 & \(1.45\times10^{-2}\) & \(5.58\times10^{-5}\) \\
\bottomrule
\end{tabularx}
\end{center}
La dernière colonne n'a été calculée qu'après fixation des sorties. Elle n'a participé ni au balayage, ni à \(G\), ni à \(K\), ni à \(Q\), ni à la transformation de Cayley, ni à la couture nonadique.\par
\Needspace{5\baselineskip}
\subsection{Convergence et préservation du registre généalogique spectral}
Après la résolution de Riemann–Weil, la fonctionnelle possède la représentation\par
\[
\mathscr I_{\rm GW}(f,g)
=
\sum_\gamma
m_\gamma
\widehat f(-\gamma)
\overline{\widehat g(-\gamma)}.
\]
Celle-ci ne sert pas à calculer les entrées : elle permet de reconnaître le spectre qui vient d'être obtenu.\par
La construction GNS s'identifie à l'espace \(L^2\) de la mesure atomique\par
\[
\nu_W=\sum_\gamma m_\gamma\delta_\gamma.
\]
Le générateur des translations agit par multiplication par \(\gamma\). Ainsi :\par
\begin{itemize}
\item \(G\) contient les produits scalaires ;
\item \(K\) insère une puissance de \(\gamma\) ;
\item \(Q\) inserta \(\gamma^2\);
\item le faisceau \(K-\gamma G\) localise les atomes ;
\item la transformation de Cayley les convertit en phases unitaires ;
\item la couture nonadique transporte ces phases dans le registre.
\end{itemize}
On n'utilise pas la résolution pour redémontrer RH avec dix matrices. On utilise la résolution déjà obtenue pour transformer la formule explicite en un algorithme d'extraction spectrale qui ne reçoit pas les zéros en entrée.\par
\Needspace{5\baselineskip}
\subsection{Convergence du champ spectral reconstruit et extension certifiable}
Le cycle de calcul est clos :\par
\[
\text{premiers+gamma+polaire}
\longrightarrow
(G,K,Q)
\longrightarrow
\widehat H_N
\longrightarrow
C_N
\longrightarrow
S_{N,9}
\longrightarrow
\{\gamma_n\}.
\]
L'interprétation littérale du cercle de clôture est également construite :\par
\[
\text{un tour nonadique complet}
\quad\Longleftrightarrow\quad
\text{une application de la phase spectrale},
\]
mais non\par
\[
\text{une phase nonadique}
\quad\Longleftrightarrow\quad
\text{un zéro}.
\]
Chaque zéro définit un caractère du registre,\par
\[
n\longmapsto\tau_\gamma^n,
\]
et dans l'expression contractive, il apparaît comme\par
\[
n\longmapsto
\left(\frac59\right)^n\tau_\gamma^n.
\]
Telle est l'orbite concrète : non une répétition de hauteurs, mais l'évolution d'une même phase spectrale à travers la mémoire des tours.\par
On ne confond pas non plus le nombre \(54\) de la branche exceptionnelle avec un nombre de zéros ou de tours. Le \(54\) d'APP–Witt–Leech–Monstre est défaut, norme et caractère gradué dans son domaine. Il peut coexister avec le cercle spectral par un produit de représentations, mais il n'a sélectionné ni normalisé aucune des hauteurs calculées ici.\par
La convergence mathématique repose sur une classe cofinale de fenêtres et de modes compacts qui constitue un cœur pour la norme du graphe du générateur GNS. Le calcul effectué fournit quatre raffinements de fenêtre et des résidus a posteriori. Une version certifiée en arithmétique d'intervalles pourrait convertir chaque décimale en un intervalle garanti, mais il ne manque plus aucune application conceptuelle ou opératorielle : ce serait un renforcement numérique d'une chaîne déjà construite de bout en bout.\par
La formule explicite et la connexion nonadique constituent la structure mathématique initiale. Les matrices \(G,K,Q\), l'opérateur de Cayley–Galerkin et l'entrelacement fini du double cercle constituent la construction numérique. Les dix hauteurs sont extraites directement de la forme arithmétique et constituent le résultat de ce calcul.\par
